\section{引言}

本文来自 Anthropic Labs 的工程师 Prithvi Rajasekaran，探讨了如何通过 \textbf{multi-agent harness} 设计来突破 Claude 在两类任务上的性能天花板：\textbf{高质量前端设计}（主观审美判断）和\textbf{长时间自主软件工程}（可验证的正确性与可用性）。作者从 GAN（生成对抗网络）中汲取灵感，设计了 generator-evaluator 分离架构，并最终扩展为 planner-generator-evaluator 三 agent 流水线，在多小时的自主编码会话中产出完整的全栈应用。

\begin{importantbox}{核心论点}
\textit{``Every component in a harness encodes an assumption about what the model can't do on its own, and those assumptions are worth stress testing, both because they may be incorrect, and because they can quickly go stale as models improve.''}

Harness 中的每个组件都编码了一个关于模型能力边界的假设。随着模型迭代，这些假设需要持续验证和更新。
\end{importantbox}

